\documentclass[11pt,a4paper]{article}
\usepackage[utf8]{inputenc}
\usepackage[T1]{fontenc}
\usepackage[margin=1in]{geometry}
\usepackage{xcolor}
\usepackage{hyperref}
\usepackage{lmodern}
\usepackage{microtype}
\definecolor{primarycolor}{RGB}{33,37,41}
\definecolor{secondarycolor}{RGB}{0,102,204}
\hypersetup{colorlinks=true,linkcolor=secondarycolor,urlcolor=secondarycolor,citecolor=secondarycolor}
\pagestyle{empty}
\setlength{\parindent}{0pt}
\setlength{\parskip}{0.75em}
\begin{document}
\textbf{\color{primarycolor} Tobias Mwangi-Hartley}

Quillstone Payments\\
Hiring Team, Customer Operations AI

\textbf{Application: Machine Learning Engineer - Customer Support AI}

Dear Hiring Team,

At Thistlebox, I have built forecasting and churn models that inform operational decisions, alongside tools that help product colleagues evaluate experiments. I am applying for the Machine Learning Engineer - Customer Support AI role to bring this applied modelling experience to your team and develop further on the production side of machine learning. Your team's work with backend engineers, data scientists and product managers to take ideas from prototype to production makes this a relevant next step.

My demand-forecasting work uses LightGBM to predict weekly recipe orders six weeks ahead. Calendar, price and menu-position features helped reduce WAPE from 23\% to 14\%, with an estimated 9\% reduction in food waste. The workflow runs as a scheduled notebook in Vertex AI Workbench, writes to BigQuery and supplies results to operations through Looker. I also built an XGBoost churn model for weekly batch scoring; validation using a 10\% holdout showed churn 2.1 percentage points lower in the treated group.

Alongside modelling, I have been strengthening the engineering around my work. Since 2025, I have written pytest unit tests for forecasting feature code, reaching approximately 60\% coverage. For a Streamlit experiment calculator used by around 25 product and growth colleagues, I wrote the application and Dockerfile, while the data platform team handles Cloud Run deployment. My experience is therefore in scheduled and batch ML workflows rather than independently deployed model APIs. An MSc in Statistics supports my approach to evaluation, and a PyTorch course project gave me initial experience fine-tuning a small image-classification model.

I would bring Python, SQL, feature engineering and careful model evaluation to your customer-support AI work, while building on my existing foundation in a machine learning engineering role. I have also delivered "Stats for PMs" three times, experience relevant to Quillstone's emphasis on sharing context openly. I would welcome a conversation about how my practical modelling experience and developing engineering skills could contribute to your team.

Kind regards,\\
Tobias Mwangi-Hartley
\end{document}